% Compile with XeLaTeX (uses macOS system fonts Charter, Avenir Next, Menlo):
%   xelatex CONTRIBUTIONS.tex
\RequirePackage{fix-cm}
\documentclass[10pt]{article}
\usepackage{fontspec}
\usepackage[letterpaper,left=0.85in,right=0.85in,top=0.58in,bottom=0.55in]{geometry}
\usepackage{xcolor}
\usepackage{booktabs}
\usepackage{tabularx}
\usepackage{array}
\usepackage[hidelinks]{hyperref}

% ---------- fonts ----------
\setmainfont{Charter}
\newfontfamily\headfont{Avenir Next}[UprightFont={* Demi Bold},BoldFont={* Bold}]
\newfontfamily\labelfont{Avenir Next}[UprightFont={* Medium},BoldFont={* Demi Bold}]
\newfontfamily\metafont{Avenir Next}[UprightFont={* Regular}]
\setmonofont{Menlo}[Scale=0.84]

% ---------- colors ----------
\definecolor{navy}{HTML}{1B365D}
\definecolor{body}{HTML}{1F1F1F}
\definecolor{slate}{HTML}{4A5868}
\definecolor{rulegray}{HTML}{B9C3CF}
\definecolor{tint}{HTML}{EEF2F7}

% ---------- base layout ----------
\makeatletter
\renewcommand\normalsize{\@setfontsize\normalsize{10}{12.7}}
\makeatother
\normalsize
\color{body}
\pagestyle{empty}
\setlength{\parindent}{0pt}
\setlength{\parskip}{0pt}
\setlength{\emergencystretch}{1.5em}

% ---------- headings ----------
\newcommand{\sectionhead}[1]{%
  \par\vspace{8pt}%
  {\headfont\fontsize{12.5}{15}\selectfont\color{navy}#1\par}%
  \vspace{2.5pt}{\color{rulegray}\hrule height 0.6pt}\vspace{4pt}}
\newcommand{\subhead}[1]{%
  \par\vspace{4.5pt}%
  {\labelfont\bfseries\fontsize{10}{12}\selectfont\color{navy}#1\par}\vspace{1.5pt}}

% ---------- lists (one level only) ----------
\newlength{\itemgap}\setlength{\itemgap}{1.5pt}
\newenvironment{points}{\begin{list}{\textcolor{navy}{\small\textbullet}}{%
  \setlength{\leftmargin}{1.05em}\setlength{\labelwidth}{0.7em}%
  \setlength{\labelsep}{0.4em}\setlength{\itemsep}{\itemgap}%
  \setlength{\parsep}{0pt}\setlength{\topsep}{1pt}\setlength{\partopsep}{0pt}}}%
  {\end{list}}
\newenvironment{numbered}{\begin{list}{\labelfont\color{navy}\arabic{enumi}.}{\usecounter{enumi}%
  \setlength{\leftmargin}{1.25em}\setlength{\labelwidth}{0.9em}%
  \setlength{\labelsep}{0.4em}\setlength{\itemsep}{\itemgap}%
  \setlength{\parsep}{0pt}\setlength{\topsep}{1pt}\setlength{\partopsep}{0pt}}}%
  {\end{list}}

% ---------- helpers ----------
\newcommand{\kstar}{\textit{k}\textsuperscript{*}}
\newcommand{\fp}[1]{{\leavevmode\ttfamily\color{slate}#1}}
\newcommand{\lbl}[1]{{\leavevmode\labelfont\fontsize{9}{11}\selectfont\color{slate}#1}}
\newcommand{\num}[1]{\textbf{#1}}
\newcolumntype{L}{>{\raggedright\arraybackslash}X}

\begin{document}

% =====================================================================
% HEADER
% =====================================================================
{\headfont\fontsize{18}{21}\selectfont\color{navy}Contributions to the GDoT Data-Sufficiency Project\par}
\vspace{4pt}
{\metafont\fontsize{9.5}{12}\selectfont\color{slate}%
{\labelfont\color{body}Kevin Li}\quad\textbar\quad April\,\textendash\,August 2026\hfill
\href{https://github.com/lzhkevin99-cell/gdot-data-sufficiency-contributions}{github.com/lzhkevin99-cell/gdot-data-sufficiency-contributions}\par}
\vspace{6pt}{\color{navy}\hrule height 1pt}\vspace{0pt}

% =====================================================================
% OVERVIEW
% =====================================================================
\sectionhead{Project Overview}

When a Georgia DOT continuous count station (CCS) misses days of data, can its monthly average daily
traffic still be trusted? The project's main paper answers this with a conformal-prediction procedure
that estimates \kstar, the minimum number of observed days needed. Alongside the paper, I carried out
two independent workstreams that use its code and data read-only:
\begin{points}
  \item \textbf{Exploratory data analysis} of the full 286-station network (2017--2024), run as an
    AI-agent pipeline that I designed.
  \item \textbf{A cluster-pooled reference-pool fallback}, after I found that the paper's pipeline
    silently drops 55\% of eligible stations: a self-directed learning experiment, validated against
    baselines and not adopted.
\end{points}

\vspace{5pt}
\noindent\colorbox{tint}{\parbox{\dimexpr\linewidth-2\fboxsep\relax}{\vspace{2.5pt}%
  {\labelfont\bfseries\fontsize{9.5}{12}\selectfont\color{navy}My role\par}\vspace{1pt}
  AI coding agents executed the code and analysis from specifications I wrote. I contributed the
  problem framing, the design of the pipeline and experiments, the specifications, and the evaluation
  of every result. Each number here was checked against the output that produced it.\vspace{2.5pt}}}

% =====================================================================
% 1. EDA
% =====================================================================
\sectionhead{1\quad Exploratory Data Analysis of the CCS Network}

\subhead{Objective}
Establish what data exists, which parts are reliable, and how stations behave and relate, before any
modeling.

\subhead{Pipeline design and methodology}
\begin{points}
  \item Divided the analysis among \textbf{seven analysis agents and one synthesis agent}, each owning
    one question type, under a 475-line specification I wrote (\fp{eda/claude.md}).
  \item Ordered the agents by data dependency. The quality audit runs before any behavioral analysis,
    so untrustworthy station-months are never treated as signal:
\end{points}
\vspace{4pt}
{\centering\labelfont\fontsize{8.75}{11}\selectfont\color{slate}%
Inventory \enspace$\rightarrow$\enspace Quality audit \enspace$\rightarrow$\enspace
[\,Profiles \,\textperiodcentered\, Trends \,\textperiodcentered\, Corridors \,\textperiodcentered\,
Anomalies, in parallel\,] \enspace$\rightarrow$\enspace Features \enspace$\rightarrow$\enspace Synthesis\par}
\vspace{4pt}
\begin{points}
  \item Built two guardrails into the specification: a mandatory Risks/Caveats section for every agent,
    and explicit rules against overclaiming (e.g.\ no MCAR/MAR labels without evidence).
\end{points}

\subhead{Key findings}
{\renewcommand{\arraystretch}{1.12}%
\begin{tabularx}{\linewidth}{@{}>{\raggedright\arraybackslash}p{1.15in}L@{}}
\lbl{Inventory} & \num{286} stations with data out of a 330-station catalog, across three data tiers
  (5-min 2024, 5-min 2022--23, hourly 2017--22); 225 (78.7\%) appear in all three.\\\addlinespace[3.5pt]
\lbl{Missingness} & \num{21,837} station-months audited, \num{91.6\%} complete ($\geq$80\% bin
  coverage). Peaks in 2017 (13.6\%) and 2022--23 ($\approx$12\%, a data-pipeline transition) show it is
  operationally driven, not random.\\\addlinespace[3.5pt]
\lbl{Outage styles} & A station-level classification not in the main paper: always-on 73.4\%,
  occasional-gap 15.7\%, seasonal-dropout 9.1\%, chronic-sparse 1.8\%.\\\addlinespace[3.5pt]
\lbl{Trust} & \num{243 of 286} stations (84.9\%) rated HIGH-trust, and 19,963 station-months (96.9\%)
  trusted. Of 391 anomaly-days, 90.5\% are explained by logged events.\\\addlinespace[3.5pt]
\lbl{Corridors} & 2,848 valid station pairs (mean $r = 0.573$); \num{27} reference-station candidates
  pass all five screening criteria.\\\addlinespace[3.5pt]
\lbl{Modeling readiness} & A 286\,$\times$\,27 station feature table and an 818,645-record clean
  hourly 2024 dataset.\\
\end{tabularx}}

\subhead{Limitations and deliverables}
\begin{points}
  \item \textit{Limitations, reported by the pipeline itself rather than hidden:} holiday-effect
    estimates return NaN for all six holidays (a station-ID join mismatch), and $k{=}4$ archetype
    clustering labels 77 of 78 stations ``Mixed'' (too few features to discriminate).
  \item \textit{Deliverables:} a 14-page final report, a 14-slide final presentation, a 19-page summary
    report, a pipeline-design document, and 36 figures.
\end{points}

% =====================================================================
% PAGE 2
% =====================================================================
\newpage
% Page 2 has less text than page 1, so its spacing is slightly looser to keep the pages balanced.
\setlength{\itemgap}{4.5pt}
\renewcommand{\sectionhead}[1]{%
  \par\vspace{15pt}%
  {\headfont\fontsize{12.5}{15}\selectfont\color{navy}#1\par}%
  \vspace{2.5pt}{\color{rulegray}\hrule height 0.6pt}\vspace{6pt}}
\renewcommand{\subhead}[1]{%
  \par\vspace{8pt}%
  {\labelfont\bfseries\fontsize{10}{12}\selectfont\color{navy}#1\par}\vspace{2.5pt}}

\sectionhead{2\quad Cluster-Pooled Reference-Pool Fallback}
{\itshape\color{slate}A self-directed learning exercise: not adopted, and never presented as a change to the paper.\par}
\vspace{1pt}

\subhead{Problem identified}
The paper's multi-site pipeline requires a reference pool of at least 19 correlated companion-months
per target station-month. Tracing the code, I found that for April 2024 only \num{84 of 185} eligible
stations meet this bar. The other \num{101 (55\%)} are dropped silently: the function returns
\texttt{None}, with no record of which stations failed or why. They skew low-volume and rural, exactly
the sites a DOT most needs an answer for.

\subhead{Hypothesis and experimental design}
\textit{Hypothesis:} stations with a similar weekly traffic pattern can supplement a too-small reference pool.
\begin{points}
  \item \textit{Profiles:} a 7-value day-of-week vector divided by the station's mean, so clusters
    reflect pattern, not road size.
  \item \textit{Clustering:} K-means over $K = 2$--$12$. Silhouette peaks at $K{=}3$ (0.406);
    Calinski--Harabasz and Davies--Bouldin favor $K{=}4$--$6$. I kept $K{=}3$ by the pre-stated rule and
    documented the disagreement rather than call it optimal.
  \item \textit{Integration:} a 0.90 completeness gate leaves 194 stations (clusters of 126 / 12 / 56).
    The station's own pool is always tried first via the paper's unchanged code; an assertion blocks
    target-data leakage.
  \item \textit{Validation:} a non-circular holdout. 20 stations with large pools ($\geq$38) supply
    ground-truth \kstar; each is shrunk below 19 and re-solved by the fallback and two simpler baselines.
\end{points}

\subhead{Results and baseline comparison}
{\renewcommand{\arraystretch}{1.12}%
\begin{tabularx}{\linewidth}{@{}>{\raggedright\arraybackslash}p{1.55in}L@{}}
\lbl{Coverage: improved} & \num{84 $\rightarrow$ 95} of 185 eligible stations (45\% $\rightarrow$
  51\%); 11 stations rescued.\\
\lbl{Accuracy: not improved} & Higher \kstar\ error than the simplest baseline (table below).\\
\end{tabularx}}

\vspace{9pt}
{\centering\fontsize{9.5}{11.5}\selectfont
\setlength{\tabcolsep}{8pt}
\begin{tabular}{@{}lrrrrr@{}}
\toprule
{\labelfont\bfseries Method} & {\labelfont\bfseries \textit{n}} & {\labelfont\bfseries MAE (days)} &
  {\labelfont\bfseries RMSE} & {\labelfont\bfseries Bias} & {\labelfont\bfseries Corr.}\\
\midrule
Cluster fallback (this work)      & 100 & 3.48 & 4.85 & $+$0.68 & 0.735\\
Baseline: same functional class   &  96 & \textbf{2.91} & \textbf{4.30} & $-$1.26 & \textbf{0.818}\\
Baseline: pool all stations       & 108 & 3.18 & 4.89 & $+$2.10 & 0.787\\
\bottomrule
\end{tabular}\par\vspace{2.5pt}
{\metafont\fontsize{8.25}{10}\selectfont\color{slate}%
\kstar\ error vs.\ ground truth, 20 holdout stations $\times$ 6 missingness mechanisms
(\textit{n} = resolved pairs). Best values in bold.}\par}
\vspace{4pt}

\subhead{Final decision and interpretation}
\begin{points}
  \item \textit{Not adopted.} The simplest baseline won. A coverage gain that cannot be shown to
    preserve accuracy is not yet a contribution; the bar is the same-functional-class baseline, not
    ``better than nothing.'' Deferred work (other algorithms, sensitivity, stability) is tagged in the run
    report.
  \item \textit{Shape normalization is load-bearing,} independent of adoption. Raw-volume clustering
    agrees with the shape-based one at only \num{ARI\,=\,0.028}: it merely splits the 18 busiest
    stations from 176 others whose weekly shapes are nearly identical.
\end{points}

% =====================================================================
% 3. DIRECTIONS
% =====================================================================
\sectionhead{3\quad Research Directions (Documented, Not Yet Tested)}
\begin{numbered}
  \item \textit{Mechanism-conditional pooling:} select companions conditional on the missingness
    mechanism under test.
  \item \textit{Outage co-occurrence as a negative signal:} reject companions whose outages coincide
    with the target's.
  \item \textit{Impute-then-conform:} fill small gaps from behavioral neighbors before the sufficiency test.
  \item \textit{Outage style as a \kstar\ stratifier:} join the EDA's outage classes to existing \kstar\
    results.
  \item \textit{Exchangeability check:} do reference months spanning the 2022--23 transition violate the
    paper's assumption?
\end{numbered}

% =====================================================================
% 4. DEMONSTRATES
% =====================================================================
\sectionhead{4\quad What This Work Demonstrates}
\begin{points}
  \item \textit{Baselines before claims:} the coverage gain looked like success until the baseline
    comparison.
  \item \textit{Reading research code closely:} tracing control flow exposed an undocumented 55\%
    station loss.
  \item \textit{Guardrails for AI-assisted analysis:} required caveats made the outputs report their own
    failures.
\end{points}

\end{document}
